\begin{lemma}[Missing colors change only at PE supporting vertices]
Let \(\mathcal T\) be a certified Tashkinov series, let \(p,q<r\) with
\(p\leq q\), and put \(i=p+1\) and \(n=q+1\).  If
\(v\in V(T_i)\) and no iteration \(j\) in the closed interval
\(i\leq j\leq n\) is a PE whose chosen supporting vertex is \(v\), then
\[
 \overline\varphi_{i-1}(v)=\overline\varphi_n(v).
\]
Here Lean index \(a\) denotes iteration \(a+1\), so the quantified
interval is exactly \(p\leq a\leq q\), not an interval shifted at either
end.
\end{lemma}
\begin{proof}
We prove the stronger running assertion that after processing every
iteration \(a+1\) with \(p\leq a\leq q\), the missing-color set at \(v\)
is unchanged from \(\overline\varphi_p(v)\).  The ordered-extension field
in each noninitial constructor of \noderef{TashkinovSeriesStep}, together
with \noderef{TreeSequenceInitialSegment}, implies inductively that
\(v\in V(T_{a+1})\) throughout the interval.

Fix one such \(a\) and eliminate the certified step supplied by
\noderef{TashkinovSeries}.  In the RE constructor the current coloring is
definitionally the previous coloring, so there is nothing to prove.  In
the SE constructor the current coloring is the selected
\((T_{a+1},D_a,\varphi_a)\)-stable coloring.  The second clause of
\noderef{StableColoring} gives equality of the two missing-color sets at
every vertex of \(T_{a+1}\), in particular at \(v\).

Consider a PE.  Let \(x\) be its supporting vertex, let \(\pi\) be the
first selected stable coloring, and let \(\rho\) be the second stable
coloring.  Stability first gives
\(\overline\pi(w)=\overline\varphi_a(w)\) and then
\(\overline\rho(w)=\overline\pi(w)\) for every
\(w\in V(T_{a+1})\).  These are direct applications of
\noderef{StableColoring}; equivalently their composition is the
transitivity supplied by \noderef{sc2} after restricting the enlarged
stored-color set to \(D_a\).

It remains to analyze the component switch.  The PE constructor in
\noderef{TashkinovSeriesStep} supplies a \noderef{TreeExit} from an
outside vertex to \(x\) whose path has no other tree vertex, with its
endpoint and support conditions given by \noderef{TreeExitPath}.  The
rule \noderef{ParallelExtensionColor} says that the PE color
\(\gamma\) is missing at \(x\), while its connecting edge supplies the
other color \(\delta\) there.  Thus \(x\) has degree one in the
two-color incidence graph.  At the outside endpoint, one of the two colors
is missing and the path supplies the other, so that endpoint also has
degree one.  By \noderef{TwoColorIncidenceDegreeAtMostTwo}, all degrees in
the component are at most two.  The usual first-branch argument therefore
shows that the supplied simple exit path exhausts the full component: an
extra component node would attach either to an internal path node, giving
degree at least three, or to an endpoint, giving degree at least two.
Consequently the full switched component meets the current tree precisely
at \(x\).

By \noderef{TwoColorComponentSwap}, switching that component changes
incident colors, and hence missing colors, at the one tree vertex \(x\)
only.  The interval hypothesis says that the present PE does not have
\(x=v\).  Thus the missing-color set at \(v\) is unchanged in the PE case
as well.

Induction over the finite ordered interval \(p,p+1,\ldots,q\) composes
these one-step equalities.  Its initial set is
\(\overline\varphi_p(v)=\overline\varphi_{i-1}(v)\), and its final set is
\(\overline\varphi_{q+1}(v)=\overline\varphi_n(v)\), which proves the
claim.
\end{proof}
